% Pursuing Stacks (cotes 134-2 to 134-8): archive cote 134-3, site page 18 (PDF page 19), typescript page 18.
% typescript. Draft compared: ps2.tex:806-864.
% AI-transcribed by Claude Opus 5.5 (claude-opus-5-5, adaptive thinking, effort high), Message Batches API, batch ps1 (msgbatch_01Rv5hCaNTtUt4JHx1yqoHPL), 2026-11-12; not yet checked by eye.
\scanpage{134-3/18}
\mspage{18}

There is another amusing interpretation of the homotopy type associa-%
ted to any preordered set K,  via the topological space whose underlying
\margin{$\times$} set is K itself,  and where the closed sets are the subsets J of K
such that  $x \hand{\in} J$,  $y \hand{\unsure{\leq|\subseteq}} x$ implies $y \hand{\in} J$.  This is a highly non-separated
topology\inserted{$\mathcal{T}$} (except when the preorder relation is the discrete one),  where
an arbitrary union of closed subsets is again closed.  I doubt it's
singular homotopy type to make much sense,  however it's  homotopy type
as a topos does,  and (possibly under mild loc\unsure{\hand{-}}al finiteness restriction\unsure{s}\hand{)}
it should be the same as the homotopy type of K just envisioned.  Thus,
a sheaf on the topological space K  \xout{defines} can be interpreted \xout{as being}
\unsure{V}ia its fibers as being just a \uline{covariant} functor
\[ K \hand{\longrightarrow} (\text{Sets}) \]
(NB  the open sets of K are just the closed sets of $K^{o}$ , namely for the
opposite order relation,  and thus every $x \hand{\in} K$ has a smallest neighbour-%
hood,  namely the set \xout{of all} $K_{\hand{\geq} x}$  of all $y \hand{\in} K$ such that $y \hand{\geq} x$),  or wh\hand{at}
amounts to the same,  a sheaf on the topos  $\hand{\widehat{K^{o}}}$  defined by the opposite
order.  Hence the derived functors of the ``sections'' functor,  when
working with abelian sheaves\inserted{on $(K,\mathcal{T})$} \xout{is just the} i.e.  the cohomology of the
topological space\inserted{$(K,\mathcal{T})$} can be interpreted in terms of the topos associa-%
ted to $K^{o}$ .  This suggests that the definition I gave of the topology
of K was awkward and maybe it is indeed (although it is the more nat\over{\unsure{y}}{u}ral
one in terms of incidence relations between open strata of $\hand{|}K\hand{|}$ ),\inserted{and} that
we should have called ``open'' the sets I called ``closed'' and vide-versa,
or equivalently,  replace K by $K^{o}$ in the definition I gave of a topology
on the set K.  But as far as homotopy types are concerned, it does'nt
make a difference,  namely the homotopy types associated to K and $K^{o}$ are
canonically isomorphic.  This can be seen most simply on the topological
realizations,  via a homeomorphism (not only a homotopism)
\[ \hand{|}K\hand{|} \hand{\simeq} \hand{|}K^{o}\hand{|} , \]
coming from the fact that the maquettes (by which I mean the combinatori%
\margin{(H.)} al model  for a triangulation,  which\inserted{H.} Cartan time ago called ``schéma sim%
plicial'' ...) of the two spaces are canonically isomorphic,  because
the ``flags'' of K and $K^{o}$ are the same.  A similar argument due to Quillen
using the nerves shows that for any category K  (not necessarily ordered\unsure{\hand{)}}
K and $K^{o}$ are homotopic,  although\inserted{\unsure{$\widehat{K}$ and $\widehat{K^{o}}$|$K$ and $K^{o}$} are} definitely not equivalent i.e. \inserted{not}
``homeomorphic'' .

To come back to the decreasing cascade of algebraic structures sui-%
table for describing homotopy types,  we could go down one more step
still,  to the category of  maquettes  (Maq) ,  namely sets\inserted{S} together
with a family\inserted{K} of finite subsets (the simplices\inserted{\unsure{S}} of \unsure{t}he set of vertices
